\section{完整物理系统与降阶动力学}

\subsection{相量约定、幅值约定与 Maxwell 方程}
全文采用 $e^{+i\omega t}$ 相量约定，并统一使用峰值复幅值。因此周期平均功率均带 $1/2$：
\begin{equation}
P_{\rm in}=\frac12\operatorname{Re}(c^H v).
\end{equation}

对给定几何 $g$，当前材料模型中的电磁 constitutive parameters 不随温度变化，故
\begin{equation}
\nabla\times\mu^{-1}\nabla\times E-\omega^2\epsilon E+i\omega\sigma E=-i\omega J_s,
\end{equation}
离散后写成
\begin{equation}
\boxed{A_{\rm em}(g)X(g)=B(g),\qquad B(g)=-i\omega S(g).}
\end{equation}
$X\in\C^{n_E\times n_p}$ 包含所有端口基激励响应。未来若加入真正影响 Maxwell constitutive law 的热状态，只需恢复 $A_{\rm em}(s_{\rm em},g)$，不改变后续端口/张量结构。

\subsection{频域电磁--瞬态热耦合的时间尺度前提}
生产模型采用 frequency--transient 分层，要求载波场建立时间远短于 thermal/envelope time scale。thermal model 接收 cycle-averaged Joule heat，而不解析载波周期内温度微振荡。若 excitation 出现与电磁瞬态同量级的快速调制、宽带/多谐波或强非线性频谱，则当前单频相量模型失效。

\subsection{离散算子与开放域结构}
Maxwell 离散保持
\begin{equation}
A_{\rm em}=C^TH_{\mu^{-1}}C-\omega^2H_\epsilon+i\omega H_\sigma+A_{\rm open},
\end{equation}
其中 $A_{\rm open}$ 为经过 Physics Gate 验证的 reciprocal/passive open-boundary contribution。Joule、端口功率与 thermal heat 必须从同一 loss-operator 体系构造。

端口 source $S_p$ 必须有明确 orientation、terminal 与 feed/return/charge-continuity 解释。开放 two-terminal spiral 不能错误地强制为闭合 divergence-free loop。source regularization 也必须绑定物理 conductor scale，不能让 mesh size 充当隐式 wire radius。

\subsection{热方程采用温升变量}
令
\begin{equation}
\theta=T-T_{\rm amb}.
\end{equation}
当前 thermal constitutive model 中 $M_T(g)$ 与 $K_T(g)$ 随 geometry 变化但不随温度变化，完整热模型为
\begin{equation}
\boxed{
M_T(g)\dot\theta+K_T(g)\theta
=Q_{\rm vol}(g,c)+Q_{\rm wire}(\theta,g,c)+F_T(g).
}
\end{equation}
理论默认生产域内
\begin{equation}
M_T(g)\succ0,\qquad K_T(g)\succ0,
\end{equation}
且 conditioning/coercivity 足够统一。若未来采用 pure Neumann 等使 $K_T$ 含 nullspace，则必须重新处理 steady-state solvability、平均温度漂移与时间积分理论。

\subsection{固定全局 thermal basis 不再作为生产假设}
对大范围平移/旋转的 coil/package，短时热响应会携带局部热点的位置自由度。若强制
\begin{equation}
\theta\approx\Phi a,\qquad \Phi\text{ 与 }g\text{ 无关},
\end{equation}
则固定线性空间必须同时记忆所有可能的热点位置，rank 会随 geometry coverage 快速增长，并可能出现 training anchors 已收敛而 held-out geometry 仍近乎正交于该 span 的情况。

因此生产理论改为 geometry-aware deterministic basis：
\begin{equation}
\boxed{
\theta(x,t;g)\approx\Phi(x;g)a(t;g).
}
\end{equation}
首版 basis generator 定义为
\begin{equation}
\boxed{
\Phi(g)=
[\Phi_{\rm bg},\;\mathcal T_{\rm tx}(g)\Psi_{\rm tx},\;\mathcal T_{\rm rx}(g)\Psi_{\rm rx}].
}
\end{equation}
其中 $\Phi_{\rm bg}$ 为固定 global/background modes，$\Psi_{\rm tx},\Psi_{\rm rx}$ 在各自 canonical local coordinates 中定义。$\mathcal T_{\rm tx/rx}(g)$ 首版只实施已知 rigid translation/rotation 与固定守恒插值。

不同 geometry 上，列的空间位置可变，但 block 顺序、列数和 mode index 的物理语义必须固定。禁止通过 geometry-dependent arbitrary eigensolver ordering 或未对齐 POD 使 mode identity 在样本间交换。

\subsection{Geometry-aware Galerkin ROM}
对一次静态 geometry query，$g$ 在整个 thermal trajectory 中固定，因此 $\Phi(g)$ 也固定。Galerkin 投影得到
\begin{equation}
\boxed{
M_r(g)\dot a
=-K_r(g)a+q_{\rm vol}(g,c)+q_{\rm wire}(a,g,c)+f_T(g),
}
\end{equation}
其中
\begin{equation}
\boxed{
M_r(g)=\Phi(g)^TM_T(g)\Phi(g),\qquad
K_r(g)=\Phi(g)^TK_T(g)\Phi(g).
}
\end{equation}
真实 $M_r,K_r$ 每个 geometry 确定性组装，不由神经网络预测。生产域内必须验证 $\Phi(g)$ 满列秩、$M_r(g),K_r(g)\succ0$ 且 conditioning 可接受。

若未来允许 $g=g(t)$，则
\begin{equation}
\dot\theta=\Phi(g)\dot a+\dot\Phi(g)a,
\end{equation}
从而 reduced dynamics 会出现
\begin{equation}
\Phi^TM_T\dot\Phi\,a
\end{equation}
等 transport terms。当前静态 geometry ROM 不允许忽略这些项后直接用于运动几何。

\subsection{Canonical thermal mode 构造}
固定 geometry 下，volume Joule heat 对 complex current $c$ 是 Hermitian quadratic map，因此 current-induced spatial heat family 的实维数至多为 $n_p^2$。source anchors 使用能够张成完整 Hermitian port space 的确定性 current combinations，而不依赖随机 current sampling。

canonical mode library 分为 background、TX-local、RX-local blocks。对 reference/global 或 local-frame thermal problems，统一使用 resolvent anchors
\begin{equation}
A_su=b,\qquad A_s=K+sM,
\end{equation}
其中 source anchor 取 $b=Q_\alpha$，initial-condition anchor 取 $b=M\theta_{0,\beta}$。Galerkin approximation 的理论相对能量误差为
\begin{equation}
\boxed{
\eta_s^2
=\frac{r^TA_s^{-1}r}{b^TA_s^{-1}b}
=\frac{\|u-\tilde u\|_{A_s}^2}{\|u\|_{A_s}^2}.
}
\end{equation}
rank 自动选择的对象是 canonical block modes，而不是把不同全局位置的局部 hotspot 逐一加入固定 basis。

\subsection{时间尺度必须与可解析 thermal physics 一致}
对 diffusion coefficient $\alpha$，时间尺度 $\tau$ 的 diffusion length 约为
\begin{equation}
\ell_d(\tau)=\sqrt{\alpha\tau}.
\end{equation}
若 $\ell_d(\tau)$ 远小于 thermal spatial discretization 可解析尺度，则 $(K+sM)^{-1}Q$ 在 $s=1/\tau$ 下主要反映离散 source placement，而非可解析热扩散。这样的极短时 anchor 不默认用于 production rank target，除非应用确实要求该时间尺度且 full-order mesh 已能解析。

首版优先覆盖可解析的短/中时尺度并加入 $s=0$ steady anchor。长时间有限时刻不需要一一对应的巨大 $\tau$ anchor；同一个稳定 reduced ODE 通过连续/链式时间积分到达任意查询 horizon，steady state 则独立求解并验证稳定性。

\subsection{初始条件语义}
如果输入 full-order 温升 $\theta_0$，默认采用当前 geometry 的 thermal-mass projection
\begin{equation}
\boxed{
a_0=
[\Phi(g)^TM_T(g)\Phi(g)]^{-1}\Phi(g)^TM_T(g)\theta_0.}
\end{equation}
并报告 initial projection error。允许的 initial-condition family 还必须通过 homogeneous resolvent/trajectory audit 覆盖。

若用户直接提供 reduced coordinates $a_0$，它只定义 ROM 内部状态；只有属于冻结/审计 initial-condition family 时才携带 full-order accuracy 声明。

\subsection{线圈温度与电阻状态}
每个线圈的 wire-temperature state 必须先由物理温度泛函定义，例如
\begin{equation}
\bar\theta_p(a,g)
=\ell_p^{(T)T}(g)\Phi(g)a,
\qquad
\ell_p^{(T)T}\mathbf 1=1,
\end{equation}
再由 constitutive law 得到 $R_p(a,g)$。wire resistance 不能把整个 reduced state 当无定义黑箱输入。

\subsection{time-domain 验证不可省略}
resolvent energy error 只是 canonical basis 构造与局部诊断工具，不是完整 trajectory certificate。冻结 held-out geometries 上必须直接比较 full thermal 与 geometry-aware ROM 的可解析短时、中间时间、长时间和 steady outputs，包括
\begin{equation}
T(x,t),\qquad T_{\min}(t),\qquad T_{\max}(t),\qquad \bar T_{\rm wire}(t),\qquad a_*.
\end{equation}

如果 held-out error 主要随 rigid translation/rotation 增大，应首先修正 local-frame transport/interpolation，而不是继续向固定 global span 堆 mode。若 rigid transport 后主要剩余误差与尺寸变化相关，才考虑 deterministic scale-aware extension。
